\begin{center}
    \chapter{State of the Art}
\end{center}
\label{chap:sota}

\section{Introduction}

The HireFlow project relies on a heterogeneous stack of
artificial-intelligence and natural-language-processing techniques to
build an end-to-end agentic recruitment assistant. This chapter
introduces the key concepts, architectures and protocols on which the
platform is built: the foundational notions of machine learning,
transformer-based language models and parameter-efficient fine-tuning,
multi-agent systems, knowledge graphs and graph-augmented retrieval,
candidate--job matching, affective computing and speech technologies,
and personalised ranking.

\section{Foundations of Artificial Intelligence and Natural Language
Processing}

This section recalls the building blocks shared by every other
component of HireFlow: the AI/ML/DL hierarchy, the three learning
regimes, transfer learning, embeddings, and the attention mechanism
that underlies the transformer architecture.

\subsection{Artificial Intelligence, Machine Learning and Deep Learning}
\emph{Artificial intelligence} (AI) is the broad scientific field
concerned with the design of computational systems that perform tasks
usually associated with human intelligence: perception, reasoning,
decision-making, language understanding and creative production.
\emph{Machine learning} (ML) gathers the methods inside this field that
improve their behaviour by inferring patterns from data rather than by
following hand-written rules. \emph{Deep learning} (DL) is the subset
of machine learning that relies on artificial neural networks with many
layers, trained end-to-end by gradient descent, and is responsible for
most of the progress in vision, speech and language over the past
decade.

\subsection{Learning Regimes}
Three learning regimes coexist in modern systems:
\begin{itemize}
    \item \emph{Supervised learning} trains a model on input--output
        pairs $(x, y)$ produced by human annotation, so that it learns
        to predict the correct $y$ for an unseen $x$.
    \item \emph{Self-supervised learning} replaces the human label
        with a pretext task derived from the data itself, such as
        predicting the next token of a sentence or the masked word
        inside a paragraph; it is the dominant paradigm for
        pre-training large language and speech models on corpora too
        vast to annotate by hand.
    \item \emph{Unsupervised learning} discovers structure in
        unlabelled data through clustering, dimensionality reduction
        or density estimation.
\end{itemize}

\subsection{Pre-training, Fine-Tuning and Transfer Learning}
\emph{Transfer learning} reuses a model trained on one task to solve a
related one. In modern natural-language processing, the canonical form
is \emph{pre-training} a large model on a generic self-supervised task
with a very large corpus, then \emph{fine-tuning} it on a small
labelled dataset that targets the application. Pre-training amortises
the cost of learning general linguistic competence, while fine-tuning
specialises the model on the downstream task at a fraction of that
cost.

\subsection{Tokenisation and Embeddings}
A \emph{tokeniser} converts a string into a sequence of integer
identifiers from a fixed vocabulary, usually through subword algorithms
such as \acs{BPE}, WordPiece or SentencePiece, which eliminate
out-of-vocabulary words. An \emph{embedding} is a fixed-size dense
vector that represents a discrete object --- a token, a sentence, a
document or an entity --- in a continuous space where geometric
proximity reflects semantic similarity. \emph{Sentence embeddings}
aggregate the contextual embeddings of a whole sentence into a single
vector, typically by mean-pooling, and the cosine similarity between
two such embeddings is the standard semantic-distance measure.

\subsection{The Attention Mechanism}
The \emph{attention} mechanism lets every position of a sequence
selectively gather information from any other position, weighted by a
learned compatibility score. In its scaled-dot-product form, three
projections of the input --- the queries $Q$, the keys $K$ and the
values $V$ --- are combined as
\begin{equation}
    \mathrm{Attention}(Q, K, V) =
        \mathrm{softmax}\!\left(\frac{QK^{\top}}{\sqrt{d_k}}\right) V,
    \label{eq:attention}
\end{equation}
\noindent which produces a context-sensitive update for each position.
\emph{Self-attention} restricts the operation to a single sequence;
\emph{cross-attention} lets one sequence attend to another; and
\emph{multi-head} attention runs several such operations in parallel
and concatenates their outputs.

\section{Large Language Models and Fine-Tuning}

This section presents the transformer architecture, the family of
large language models built on top of it, and the parameter-efficient
methods that adapt them to specialised tasks without re-training the
billions of base weights.

\subsection{The Transformer Architecture}
Modern language models almost exclusively rely on the
\textbf{transformer} architecture introduced by Vaswani \emph{et~al.}
in 2017~\cite{vaswani2017attention}. The transformer replaces sequential
recurrence with self-attention, which lifts a long-standing parallelism
bottleneck and lets the model capture dependencies between any two
positions in constant computational depth. Figure~\ref{fig:transformer}
shows the canonical encoder--decoder architecture: the encoder ingests
the source sequence and produces a contextualised representation
through stacked layers of multi-head self-attention and feed-forward
networks; the decoder generates the output sequence autoregressively,
combining masked self-attention with cross-attention over the encoder
output. Stacking identical blocks yields the encoder, the decoder, or,
in the most common contemporary configuration, a \emph{decoder-only}
stack specialised for next-token prediction.

\begin{figure}[H]
    \centering
    \includegraphics[width=0.5\textwidth]{transformer-arch.png}
    \caption{Transformer encoder--decoder architecture.}
    \label{fig:transformer}
\end{figure}

\subsection{Large Language Models}
A \emph{large language model} (LLM) is a transformer-based decoder-only
network with billions of parameters, pre-trained self-supervised on a
multi-trillion-token corpus to predict the next token, then
instruction-tuned to follow human prompts. Open-weight families such as
\textbf{Llama} and \textbf{Qwen} make a class of techniques ---
inspection, fine-tuning, on-premises deployment --- accessible that
closed commercial models do not.

\subsection{Parameter-Efficient Fine-Tuning}
Adapting an open-weight LLM to a specialised task no longer requires
full fine-tuning of the base model. The \emph{parameter-efficient
fine-tuning} (PEFT) family freezes the base weights and trains a small
number of additional parameters that capture the task-specific
adjustment. The dominant approach is \textbf{LoRA} (Low-Rank
Adaptation)~\cite{hu2022lora}, which adds to each linear layer a
low-rank update
\begin{equation}
    \Delta W = B A,\qquad
    A \in \mathbb{R}^{r \times d},\
    B \in \mathbb{R}^{d \times r},\
    r \ll d,
    \label{eq:lora}
\end{equation}
\noindent shrinking the trainable-parameter count by 100 to
$1{,}000\times$ with negligible quality loss. \acs{QLoRA} extends the
recipe by quantising the frozen base weights to four-bit precision,
which lets a multi-billion-parameter model be fine-tuned on a single
consumer-class \acs{GPU}.

\section{Multi-Agent Systems and Agent Protocols}

This section defines what an agent is, presents the ReAct loop on
which most modern agent implementations are based, surveys the
families of multi-agent orchestration frameworks, and closes with the
open communication protocols that connect them.

\subsection{Agent Definition}
An \emph{agent} is an autonomous software entity built around a
language model that, given a goal expressed in natural language,
decides on its own which actions to take, executes them through external
\emph{tools} and observes the results before choosing the next step.
Three properties distinguish this paradigm from the prompt-and-response
one. First, the agent maintains an internal \emph{state} across several
steps, which lets it pursue a multi-step plan. Second, it has access
to a set of tools --- search engines, databases, code interpreters,
custom \acs{API}s --- that ground its reasoning in the external world
and allow side effects beyond text generation. Third, it autonomously
decides when the task is complete and returns control.

\subsection{The ReAct Framework}
The \textbf{ReAct} pattern (\emph{Reason} $+$ \emph{Act}), introduced by
Yao \emph{et~al.} in 2022~\cite{yao2023react}, is the canonical
instantiation of the agent loop. As sketched in Figure~\ref{fig:react},
the model alternates between three kinds of step interleaved in a
single autoregressive trace: \emph{Thought}, in which it verbalises its
reasoning in natural language; \emph{Action}, in which it emits a
structured tool call parsed from its output; and \emph{Observation},
in which the result of the tool call is fed back into the prompt before
the next thought. The loop continues until the agent decides to commit
to a final answer. Compared with chain-of-thought reasoning alone, the
explicit observation channel grounds the chain in factual evidence and
substantially reduces hallucination on tasks that require external
information.

\begin{figure}[H]
    \centering
    \begin{tikzpicture}[
        font=\small,
        every node/.style={align=center},
        block/.style={rectangle, rounded corners, draw=black!70, thick,
            minimum height=0.8cm, minimum width=2.6cm,
            inner sep=4pt, font=\small},
        query/.style={block, fill=gray!15},
        thought/.style={block, fill=yellow!22},
        action/.style={block, fill=blue!18},
        obs/.style={block, fill=green!18},
        ans/.style={block, fill=red!22, font=\small\bfseries},
        a/.style={-{Latex[length=2mm]}, thick}
    ]
        \node[query]                       (q)  {User query};
        \node[thought, below=0.55cm of q]  (t)  {Thought};
        \node[action,  below=0.55cm of t]  (ac) {Action (tool call)};
        \node[obs,     below=0.55cm of ac] (o)  {Observation};
        \node[ans,     below=0.7cm of o]   (an) {Final answer};

        \draw[a] (q)  -- (t);
        \draw[a] (t)  -- (ac);
        \draw[a] (ac) -- (o);
        \draw[a, dashed] (o) -- (an);
        \draw[a] (o.east) to[bend right=80] node[right, font=\scriptsize\itshape, xshift=2pt]{loop until done} (t.east);
    \end{tikzpicture}
    \caption{ReAct thought--action--observation loop.}
    \label{fig:react}
\end{figure}

\subsection{Multi-Agent Orchestration}
A single agent rarely covers the breadth of a real product. Three
families of multi-agent frameworks coexist:
\begin{itemize}
    \item \emph{Crew-based} frameworks structure the collaboration as
        a team with explicit roles and a sequential or hierarchical
        workflow.
    \item \emph{Conversational} frameworks model agents as
        participants whose interactions are themselves driven by a
        language model.
    \item \emph{State-machine} frameworks express the multi-agent
        process as a typed graph whose nodes are agent invocations
        and whose edges are conditional transitions, which makes the
        dialogue auditable and the branching deterministic.
\end{itemize}

\subsection{Agent Communication Contracts}
Two open specifications dominate the discussion:
\begin{itemize}
    \item The \emph{Model Context Protocol} (\acs{MCP}), introduced
        in late~2024, describes how a host application exposes tools,
        resources and prompts to a language-model-driven client.
    \item The \emph{Agent-to-Agent} (\acs{A2A}) protocol, proposed in
        2025, governs how peer agents discover each other through
        \emph{AgentCards} and exchange tasks and task results,
        typically over HTTP and JSON-RPC.
\end{itemize}
\noindent Both protocols decouple agent identity from transport, which
makes the message contract an independent design artefact.

\section{Knowledge Graphs and Graph-Augmented Retrieval}

This section defines knowledge graphs, recalls the standard
retrieval-augmented-generation recipe and presents Graph RAG, the two
grounding mechanisms HireFlow uses to mitigate hallucination.

\subsection{Knowledge Graphs}
A \emph{knowledge graph} (\acs{KG}) represents domain knowledge as a
labelled directed graph in which nodes are entities and edges are typed
relationships, traditionally as subject--predicate--object triples
inherited from \acs{RDF}. Knowledge graphs constrain the answers of a
generative model to facts present in the graph and expose multi-hop
relationships that a flat document index cannot represent.

\subsection{Retrieval-Augmented Generation}
\emph{Retrieval-Augmented Generation} (\acs{RAG}) injects external
evidence into the prompt of a language model at inference time. In its simplest
form, the corpus is split into chunks, each chunk is embedded with a
sentence encoder and indexed in a vector store; at query time, the
top-$k$ chunks are retrieved by cosine similarity and concatenated to
the prompt. Naive RAG performs well on short factual lookups but
degrades on multi-hop questions and on entity-centric queries.

\subsection{Graph RAG}
\textbf{Graph RAG}~\cite{edge2024graphrag} replaces the flat chunk
index with a graph of entities and their relationships extracted from
the corpus by a language model and clustered into communities.
Retrieval becomes a graph traversal that returns a coherent subgraph,
possibly augmented with hierarchical community summaries, which
measurably outperforms naive RAG on global questions that require
integrating information across many documents.

\section{Candidate--Job Matching}

This section traces the progression of candidate--job matching from
bag-of-words retrieval to dense semantic similarity, topic-aware
hierarchical attention, graph neural networks and constraint-aware
shortlisting.

\subsection{Sentence Embeddings and Semantic Similarity}
The first generation of automated candidate screening relied on
bag-of-words retrieval (TF-IDF, BM25), which scores a resume against a
job description by counting overlapping terms and is therefore blind to
synonymy and context. \textbf{Sentence-BERT}~\cite{reimers2019sbert}
replaces lexical overlap with semantic similarity through a twin-encoder
fine-tuning of \acs{BERT} that produces fixed-size sentence vectors
directly comparable with cosine similarity, and the family has since
produced compact and multilingual variants used at scale.

\subsection{Topic-Aware Person--Job Fit (TAPJFNN)}
The \emph{person--job fit} literature frames the matching problem as a
binary classification: given a job description and a resume, predict
whether the candidate is a good fit. \textbf{TAPJFNN}
(\emph{Topic-Aware Person--Job Fit Neural
Network})~\cite{qin2020tapjfnn} is the state-of-the-art representative
of the family. It combines a four-level hierarchical attention with a
topic-aware mechanism in which a corpus-level latent topic model
modulates the higher attention weights so that the network focuses on
the sentences that align with the dominant topic of the job.
Figure~\ref{fig:tapjfnn} sketches the four-level hierarchy applied
symmetrically to the job and to the resume.

\begin{figure}[H]
    \centering
    \begin{tikzpicture}[
        font=\small,
        every node/.style={align=center},
        word/.style={rectangle, rounded corners=2pt, draw=black!50,
            fill=gray!10, minimum height=0.45cm, minimum width=0.7cm,
            inner sep=2pt, font=\scriptsize},
        sent/.style={rectangle, rounded corners=3pt, draw=black!60,
            fill=blue!10, minimum height=0.5cm, minimum width=1.4cm,
            inner sep=2pt, font=\scriptsize},
        sec/.style={rectangle, rounded corners=3pt, draw=black!60,
            fill=green!12, minimum height=0.55cm, minimum width=1.5cm,
            inner sep=2pt, font=\scriptsize},
        doc/.style={rectangle, rounded corners=3pt, draw=black!70,
            thick, fill=orange!18, minimum height=0.65cm,
            minimum width=2.0cm, inner sep=2pt, font=\scriptsize\bfseries},
        match/.style={rectangle, rounded corners=4pt, draw=black!80,
            thick, fill=red!15, minimum height=0.85cm,
            minimum width=2.5cm, inner sep=3pt, font=\small\bfseries},
        topic/.style={ellipse, draw=black!50, fill=yellow!22,
            minimum width=2.6cm, minimum height=0.75cm, font=\scriptsize},
        attn/.style={midway, font=\scriptsize\itshape, fill=white,
            inner sep=1pt},
        a/.style={-{Latex[length=1.6mm]}, thin}
    ]
        \node[word] (jw1) at (-6.6, 0)   {w};
        \node[word] (jw2) at (-5.9, 0)   {w};
        \node[word] (jw3) at (-5.2, 0)   {w};
        \node[word] (jw4) at (-4.5, 0)   {w};
        \node[sent] (js1) at (-6.25, 1.0) {sent};
        \node[sent] (js2) at (-4.85, 1.0) {sent};
        \node[sec]  (jc1) at (-6.25, 2.0) {ability};
        \node[sec]  (jc2) at (-4.85, 2.0) {ability};
        \node[doc]  (jd)  at (-5.55, 3.1) {Job vector};

        \node[word] (rw1) at ( 4.5, 0)   {w};
        \node[word] (rw2) at ( 5.2, 0)   {w};
        \node[word] (rw3) at ( 5.9, 0)   {w};
        \node[word] (rw4) at ( 6.6, 0)   {w};
        \node[sent] (rs1) at ( 4.85, 1.0) {sent};
        \node[sent] (rs2) at ( 6.25, 1.0) {sent};
        \node[sec]  (rc1) at ( 4.85, 2.0) {exp.};
        \node[sec]  (rc2) at ( 6.25, 2.0) {exp.};
        \node[doc]  (rd)  at ( 5.55, 3.1) {Resume vector};

        \node[topic] (tp) at (0, 1.5) {LDA topic\\distribution};

        \node[match] (mm) at (0, 3.7) {Fit score};
        \node[font=\scriptsize\itshape] at (0, 4.5) {head $\delta$};

        \draw[a] (jw1) -- (js1);
        \draw[a] (jw2) -- (js1);
        \draw[a] (jw3) -- (js2);
        \draw[a] (jw4) -- (js2);
        \draw[a] (js1) -- (jc1) node[attn, near end]{$\alpha$};
        \draw[a] (js2) -- (jc2);
        \draw[a] (jc1) -- (jd) node[attn, near end]{$\beta$};
        \draw[a] (jc2) -- (jd) node[attn, near end]{$\gamma$};

        \draw[a] (rw1) -- (rs1);
        \draw[a] (rw2) -- (rs1);
        \draw[a] (rw3) -- (rs2);
        \draw[a] (rw4) -- (rs2);
        \draw[a] (rs1) -- (rc1) node[attn, near end]{$\alpha$};
        \draw[a] (rs2) -- (rc2);
        \draw[a] (rc1) -- (rd) node[attn, near end]{$\beta$};
        \draw[a] (rc2) -- (rd) node[attn, near end]{$\gamma$};

        \draw[a, dashed, gray!70] (tp) -- (jc1);
        \draw[a, dashed, gray!70] (tp) -- (jc2);
        \draw[a, dashed, gray!70] (tp) -- (rc1);
        \draw[a, dashed, gray!70] (tp) -- (rc2);

        \draw[a, thick] (jd) -- (mm);
        \draw[a, thick] (rd) -- (mm);
    \end{tikzpicture}
    \caption{TAPJFNN four-level topic-aware hierarchical attention.}
    \label{fig:tapjfnn}
\end{figure}

\subsection{Latent Dirichlet Allocation}
The topic distributions used by TAPJFNN are produced offline by
\emph{Latent Dirichlet Allocation}~\cite{blei2003lda}, the canonical
generative topic model. LDA assumes that each document is a mixture of
a small number of latent topics and that each topic is in turn a
distribution over words; the two distributions are inferred jointly
through approximate inference.

\subsection{Graph Neural Networks}
A \emph{graph neural network} (GNN) generalises the convolution
operator from regular grids to arbitrary graphs. Each node holds a
feature vector, and the network repeatedly updates that vector by
aggregating messages from its neighbours, followed by a transformation
shared across the graph; after $L$ such layers, the representation of
a node summarises an $L$-hop neighbourhood. Figure~\ref{fig:gnn}
sketches one round of message passing on a bipartite candidate--job
graph. Three operator families dominate: \acs{GCN} uses a fixed symmetric
normalisation, \acs{GAT} replaces the fixed weights with learned
attention coefficients per edge, and \acs{GIN} uses an injective
aggregation that maximises graph-classification power.

\begin{figure}[H]
    \centering
    \begin{tikzpicture}[
        font=\small,
        every node/.style={align=center},
        cnode/.style={circle, draw=black!70, thick, fill=blue!18,
            minimum size=0.75cm, inner sep=0pt, font=\scriptsize\bfseries},
        jnode/.style={circle, draw=black!70, thick, fill=orange!22,
            minimum size=0.75cm, inner sep=0pt, font=\scriptsize\bfseries},
        msg/.style={-{Latex[length=1.7mm]}, thick, blue!60},
        plain/.style={gray!55, thick}
    ]
        \node[cnode] (c1) at (-3.5,  1.6) {C$_1$};
        \node[cnode] (c2) at (-3.5,  0.0) {C$_2$};
        \node[cnode] (c3) at (-3.5, -1.6) {C$_3$};
        \node[jnode] (j1) at ( 1.5,  1.6) {J$_1$};
        \node[jnode] (j2) at ( 1.5,  0.0) {J$_2$};
        \node[jnode] (j3) at ( 1.5, -1.6) {J$_3$};

        \draw[plain] (c1) -- (j1);
        \draw[plain] (c1) -- (j2);
        \draw[plain] (c2) -- (j1);
        \draw[plain] (c2) -- (j2);
        \draw[plain] (c2) -- (j3);
        \draw[plain] (c3) -- (j2);
        \draw[plain] (c3) -- (j3);

        \draw[msg] (c1) -- (j2);
        \draw[msg] (c2) -- (j2);
        \draw[msg] (c3) -- (j2);

        \node[rectangle, rounded corners, draw=black!70, thick,
            fill=green!12, minimum width=4.4cm, minimum height=0.8cm,
            font=\scriptsize, align=center] (upd)
            at (5.4, 0) {AGGREGATE\\$+$ UPDATE};
        \draw[msg] (j2) -- (upd);

        \node[font=\scriptsize\itshape, anchor=north]
            at (j2.south) [yshift=-0.6cm] {focus node};
        \node[font=\scriptsize, anchor=west, gray!60]
            at (-3.5, -2.4) {Candidates};
        \node[font=\scriptsize, anchor=east, gray!60]
            at ( 1.5, -2.4) {Jobs};
    \end{tikzpicture}
    \caption{GNN message passing on a bipartite candidate--job graph.}
    \label{fig:gnn}
\end{figure}

\subsection{Genetic Algorithms for Constrained Shortlisting}
A pure ranker returns a list of candidates ordered by predicted fit,
but real shortlists must respect a set of operational constraints ---
a budget on the number of interviews, a diversity target, role quotas
and minimum-fit thresholds --- that do not factorise across candidates.
\emph{Genetic algorithms} are a family of biology-inspired
metaheuristics for combinatorial optimisation. A population of
candidate solutions is iteratively transformed through three operators:
\begin{itemize}
    \item \emph{Selection} retains the highest-fitness individuals
        for the next generation.
    \item \emph{Crossover} combines two parent solutions into new
        offspring.
    \item \emph{Mutation} applies small random changes to maintain
        diversity.
\end{itemize}
\noindent Constraints are typically enforced through a \emph{death
penalty} that assigns zero fitness to infeasible solutions.
Algorithm~\ref{alg:ga-shortlist} sketches the canonical loop applied to
the shortlisting problem.

\begin{algorithm}[H]
\caption{Genetic algorithm for constrained shortlisting}
\label{alg:ga-shortlist}
\begin{algorithmic}[1]
\Require Candidate pool $C$, fitness $f$, constraints $\Phi$, population size $N$, generations $T$
\Ensure Best feasible shortlist $s^{\star}$
\State Initialise population $P_0 \gets \{s_1, \dots, s_N\}$ at random
\For{$t = 0$ \textbf{to} $T-1$}
    \ForAll{$s \in P_t$}
        \State $\hat{f}(s) \gets f(s) \cdot \mathbb{1}[\Phi(s) \text{ feasible}]$ \Comment{death penalty}
    \EndFor
    \State $P^{\text{sel}} \gets \textsc{Select}(P_t, \hat{f})$ \Comment{e.g.\ tournament or roulette}
    \State $O \gets \textsc{Crossover}(P^{\text{sel}})$
    \State $O \gets \textsc{Mutate}(O)$
    \State $P_{t+1} \gets$ top-$N$ individuals of $P_t \cup O$ by $\hat{f}$
\EndFor
\State \Return $\arg\max_{s \in P_T} \hat{f}(s)$
\end{algorithmic}
\end{algorithm}

\section{Affective Computing and Speech Technologies}

\subsection{Models of Emotion}
The affective-computing community models emotion in two complementary
ways:
\begin{itemize}
    \item A \emph{discrete} set of basic categories (happiness,
        sadness, fear, anger, surprise, disgust, contempt, neutral),
        traditionally attributed to Ekman.
    \item A \emph{continuous} space whose two dominant axes are
        \emph{valence} (pleasantness) and \emph{arousal} (activation),
        known as Russell's circumplex.
\end{itemize}
\noindent The two representations coexist in production
interview-analysis systems.

\subsection{Facial Action Units and Landmarks}
The \acs{FACS} decomposes any facial expression into a small set of
\emph{Action Units} (\acsp{AU}) corresponding to individually
contractible muscles (AU4 brow lowerer, AU12 lip corner puller). AUs
are typically estimated from facial landmarks --- several hundred
three-dimensional points predicted on the face at thirty frames per
second --- rather than from raw pixels. Pose-estimation and
hand-tracking models complete the perception channel by capturing body
language and gesture.

\subsection{Self-Supervised Speech Encoders}
The dominant audio backbone today is \textbf{wav2vec~2.0}
~\cite{baevski2020wav2vec}, a self-supervised model that learns
speech representations directly from the raw audio waveform. A small
downstream head transforms its contextual embedding into a
task-specific output: transcript, emotion class, valence or arousal.

\subsection{ASR, TTS and Voice-Activity Detection}
The interview module relies on three complementary speech components:
\begin{itemize}
    \item \acs{ASR} replaces the historical \acs{HMM} pipeline with
        end-to-end encoder--decoder transformers trained on weakly
        supervised audio.
    \item Neural \acs{TTS} has moved from autoregressive seq-to-seq
        systems (Tacotron) to non-autoregressive flow-based models
        such as VITS that run faster than real time on commodity
        \acsp{CPU}.
    \item \acs{VAD} decides whether speech is present in a short audio
        window; modern neural VAD replaces the WebRTC filter with a
        small recurrent network.
\end{itemize}

\section{Personalised Ranking}

A single global ranker is a poor approximation of any one recruiter's
preferences, since two recruiters weight cultural fit, technical depth
and communication style differently. The most pragmatic personalisation
model for a high-stakes hiring use case is \emph{logistic regression}:
its coefficients are directly readable as feature weights and its
calibrated probability output admits a natural confidence
interpretation, which matters when the system operates under
regulatory constraints.

\section{Conclusion}

This chapter laid the conceptual foundations of the work and surveyed
the families of techniques on which the rest of the report builds. The
next chapters detail, sprint by sprint, the design and implementation
of the system that integrates them.
